% !TEX root = ../libretto.tex
% Sezione: Canoni (Taizé e altri)
\momento{Canoni}
\begin{songs}{titoli}

% Fonte: testo liturgico tradizionale dei canti di Taizé, scritto a memoria; verificato dall'utente.
\beginsong{Adoramus te, Christe}[by={Taizé}]
\beginverse*
Adoramus te, Christe, benedicimus tibi,
quia per crucem tuam redemisti mundum.
\endverse
\endsong

% Fonte: LibrettoAccordi.pdf (Canti Coro), O3.
% Nell'elenco: Canone.
\beginsong{Audite, beati}[by={F. Volpintesta}]
\beginverse*
Audite, beati, verbum Dei!
\endverse
\endsong

% Fonte: testo fornito dall'utente (solo ritornello).
% Non presente nell'elenco. Messo tra i canoni su indicazione dell'utente.
\beginsong{Awesome God}[by={R. Mullins}]
\beginverse*
Our God is an awesome God,
he reigns from heaven above
with wisdom, power and love:
our God is an awesome God.
\endverse
\endsong

% Fonte: testo liturgico tradizionale dei canti di Taizé, scritto a memoria; verificato dall'utente.
\beginsong{Bonum est confidere}[by={Taizé}]
\beginverse*
Bonum est confidere in Domino,
bonum sperare in Domino.
\endverse
\endsong

% Fonte: testo liturgico tradizionale dei canti di Taizé, scritto a memoria; verificato dall'utente.
\beginsong{Christe, lux mundi}[by={Taizé}]
\beginverse*
Christe, lux mundi,
qui sequitur te
habebit lumen vitae, lumen vitae.
\endverse
\endsong

% Fonte: testo liturgico tradizionale dei canti di Taizé, scritto a memoria; verificato dall'utente.
\beginsong{Confitemini Domino}[by={Taizé}]
\beginverse*
Confitemini Domino, quoniam bonus.
Confitemini Domino, alleluia.
\endverse
\endsong

% Fonte: testo liturgico tradizionale dei canti di Taizé, scritto a memoria; verificato dall'utente.
\beginsong{Crucem tuam}[by={Taizé}]
\beginverse*
Crucem tuam adoramus, Domine,
resurrectionem tuam laudamus, Domine.
Laudamus et glorificamus.
Resurrectionem tuam laudamus, Domine.
\endverse
\endsong

% Fonte: testo fornito dall'utente. Non presente nell'elenco.
\beginsong{Evangelizare pauperibus}[by={S. Zampa}]
\beginverse*
Evangelizare pauperibus
misit me, misit me.
Evangelizare pauperibus.
\endverse
\endsong

% Fonte: testo liturgico tradizionale dei canti di Taizé, scritto a memoria; verificato dall'utente.
\beginsong{Gloria, gloria in excelsis Deo}[by={Taizé}]
\beginverse*
Gloria, gloria in excelsis Deo!
Gloria, gloria, alleluia, alleluia!
\endverse
\endsong

% Fonte: testo liturgico tradizionale dei canti di Taizé, scritto a memoria; verificato dall'utente.
\beginsong{Il Signore è la mia forza}[by={Taizé}]
\beginverse*
Il Signore è la mia forza
e io spero in lui.
Il Signor è il Salvator.
In lui confido, non ho timor,
in lui confido, non ho timor.
\endverse
\endsong

% Fonte: testo liturgico tradizionale dei canti di Taizé, scritto a memoria; verificato dall'utente.
\beginsong{In manus tuas, Pater}[by={Taizé}]
\beginverse*
In manus tuas, Pater,
commendo spiritum meum.
\endverse
\endsong

% Fonte: testo liturgico tradizionale dei canti di Taizé, scritto a memoria; verificato dall'utente.
\beginsong{Jubilate Deo}[by={Taizé}]
\beginverse*
Jubilate Deo omnis terra.
Servite Domino in laetitia.
Alleluia, alleluia, in laetitia.
Alleluia, alleluia, in laetitia.
\endverse
\endsong

% Fonte: testo liturgico tradizionale dei canti di Taizé, scritto a memoria; verificato dall'utente.
\beginsong{Laudate Dominum}[by={Taizé}]
\beginverse*
Laudate Dominum, laudate Dominum,
omnes gentes, alleluia.
\endverse
\endsong

% Fonte: testo fornito dall'utente. Non presente nell'elenco. Autore indicato dall'utente.
\beginsong{Let it rain}[by={M. W. Smith}]
\beginverse*
Let it rain, let it rain,
open the floodgates of Heaven.
\endverse
\endsong

% Fonte: testo liturgico tradizionale dei canti di Taizé, scritto a memoria; verificato dall'utente.
\beginsong{Magnificat}[by={Taizé}]
\beginverse*
Magnificat, magnificat,
magnificat anima mea Dominum.
Magnificat, magnificat,
magnificat anima mea.
\endverse
\endsong

% Fonte: testo liturgico tradizionale dei canti di Taizé, scritto a memoria; verificato dall'utente.
\beginsong{Misericordias Domini}[by={Taizé}]
\beginverse*
Misericordias Domini
in aeternum cantabo.
\endverse
\endsong

% Fonte: testo liturgico tradizionale dei canti di Taizé, scritto a memoria; verificato dall'utente.
% Nell'elenco: "Nada te turbe/Nulla ti turbi", Canone. Testo di s. Teresa d'Avila.
\beginsong{Nada te turbe}[by={Taizé}]
\beginverse*
Nada te turbe, nada te espante:
quien a Dios tiene nada le falta.
Nada te turbe, nada te espante:
solo Dios basta.
\endverse
\endsong

% Fonte: testo liturgico tradizionale dei canti di Taizé, scritto a memoria; verificato dall'utente.
% Versione italiana di "Nada te turbe".
\beginsong{Nulla ti turbi}[by={Taizé}]
\beginverse*
Nulla ti turbi, nulla ti spaventi:
chi ha Dio niente gli manca.
Nulla ti turbi, nulla ti spaventi:
solo Dio basta.
\endverse
\endsong

% Fonte: testo liturgico tradizionale dei canti di Taizé, scritto a memoria; verificato dall'utente.
\beginsong{Oculi nostri}[by={Taizé}]
\beginverse*
Oculi nostri ad Dominum Jesum,
oculi nostri ad Dominum nostrum.
\endverse
\endsong

% Fonte: testo fornito dall'utente.
\beginsong{Salvator mundi}[by={Taizé}]
\beginverse*
Salvator mundi, salva nos,
Salvator mundi, salva nos.
Salva nos, salva nos,
Salvator mundi, salva nos.
\endverse
\endsong

% Fonte: testo liturgico tradizionale dei canti di Taizé, scritto a memoria; verificato dall'utente.
\beginsong{Ubi caritas}[by={Taizé}]
\beginverse*
Ubi caritas et amor,
ubi caritas Deus ibi est.
\endverse
\endsong

% Fonte: testo liturgico tradizionale dei canti di Taizé, scritto a memoria; verificato dall'utente.
\beginsong{Veni Sancte Spiritus (Tui amoris ignem)}[by={Taizé}]
\beginverse*
Veni Sancte Spiritus,
tui amoris ignem accende.
Veni Sancte Spiritus,
veni Sancte Spiritus.
\endverse
\endsong

\end{songs}
